\begin{ouverture}[calcul]{8}{Lettres, formules et motifs}{Partie 2 · Algèbre et grandeurs}
{Variable · Substitution · Régularités · Égalités}
\objectif{Produire une formule et la calculer.}
\objectif{Décrire un motif évolutif.}
\objectif{Tester une égalité pour une valeur donnée.}
\prerequis{Calculs et carrés.}
\end{ouverture}
\pagestyle{otMain}\essentiel\fontsize{11}{15}\selectfont
\phantomsection\label{cours-8-1}
\essSection{1}{Une lettre pour plusieurs valeurs}
\begin{definition}Une lettre peut représenter un nombre qui varie. Le double de $x$ s’écrit $2x$, son successeur $x+1$, son carré $x^2$. $2x$ signifie $2\times x$ ; $x^2$ signifie $x\times x$.\end{definition}
\begin{exemple}Pour un rectangle de longueur $L$ et de largeur $\ell$, $P=2(L+\ell)$ et $A=L\ell$. Si $L=7$ et $\ell=3$, alors $P=20$ unités et $A=21$ unités carrées.\end{exemple}
\label{fig-cours-8-1-cours-motif}\figureGeo{cours-motif}

\begin{savoirfaire}[Voir la structure]Une rangée de $n$ carrés accolés utilise quatre allumettes pour le premier carré, puis trois pour chaque carré ajouté. Le nombre total est $4+3(n-1)$, soit $3n+1$. Une formule décrit toutes les étapes.\end{savoirfaire}
\clearpage
\phantomsection\label{cours-8-2}
\essSection{2}{Substituer et tester}
\begin{savoirfaire}[Remplacer chaque occurrence]Pour calculer $3x^2+2x$ en $x=2$, écrire $3\times2^2+2\times2=12+4=16$. Dans une substitution, garder les parenthèses et les priorités.\end{savoirfaire}
\begin{definition}Tester une égalité pour une valeur consiste à calculer séparément ses deux membres. Pour $x=4$, $2x+1=9$ est vraie ; pour $x=3$, elle est fausse car $7\ne9$.\end{definition}
\begin{attention}Une égalité vérifiée pour une valeur n’est pas forcément vraie pour toutes les valeurs. Une formule doit aussi respecter les unités : additionner une longueur à une aire n’a pas de sens.\end{attention}
\begin{savoirfaire}[Explorer avec un tableur]Dans un tableur, si la valeur de $x$ est en A2, la formule \texttt{=3*A2+1} calcule $3x+1$. Recopier la formule vers le bas permet de tester plusieurs valeurs et de formuler une conjecture ; cela ne constitue pas une preuve générale.\end{savoirfaire}

\clearpage\rubrique{Découverte}
\begin{blocLivre}{ex-8-1}\begin{exercice}[1]{Double et successeur}
Écrire le double de $x$ et son successeur.
\end{exercice}\end{blocLivre}
\begin{blocLivre}{ex-8-2}\begin{exercice}[1]{Lire un produit}
Écrire sans signe $\times$ : $5\times a$ et $b\times b$.
\end{exercice}\end{blocLivre}
\begin{blocLivre}{ex-8-3}\begin{exercice}[1]{Calculer une formule}
Calculer $4x+3$ pour $x=2$.
\end{exercice}\end{blocLivre}
\begin{blocLivre}{ex-8-4}\begin{exercice}[1]{Une aire}
Un rectangle mesure $x$ cm sur 5 cm. Donner son aire pour $x=7$.
\label{fig-ex-8-4}\figureGeo{lit-4}
\end{exercice}\end{blocLivre}
\begin{blocLivre}{ex-8-5}\begin{exercice}[1]{Tester une égalité}
$3x+2=14$ est-elle vraie pour $x=4$ ?
\end{exercice}\end{blocLivre}

\rubrique{Entraînement}
\begin{blocLivre}{ex-8-6}\begin{exercice}[2]{Une valeur décimale}
Calculer $2t^2+1$ pour $t=1{,}5$.
\end{exercice}\end{blocLivre}
\begin{blocLivre}{ex-8-7}\begin{exercice}[2]{Une somme ou un produit}
Dire si $3x+5$ et $3(x+5)$ sont des sommes ou des produits. Calculer pour $x=2$.
\end{exercice}\end{blocLivre}
\begin{blocLivre}{ex-8-8}\begin{exercice}[2]{Motif d’allumettes}
Pour la rangée de carrés, combien faut-il d’allumettes aux étapes 1, 2, 5 et 20 ?
\label{fig-ex-8-8}\figureGeo{lit-8}
\end{exercice}\end{blocLivre}
\begin{blocLivre}{ex-8-9}\begin{exercice}[2]{Un programme verbal}
On choisit un nombre, on ajoute 4 puis on multiplie le résultat par 3. Écrire la formule et la calculer pour 6.
\end{exercice}\end{blocLivre}
\begin{blocLivre}{ex-8-10}\begin{exercice}[2]{Une formule de tableur}
La colonne A contient 0, 2 et 5 à partir de A2. On écrit \texttt{=4*A2+3} en B2 puis on recopie. Donner les trois résultats.
\end{exercice}\end{blocLivre}

\rubrique{Problèmes}
\begin{blocLivre}{ex-8-11}\begin{exercice}[3]{Une bordure de points}
À l’étape $n$, on dessine un carré ayant $n$ points sur chaque côté, sommets compris, avec $n\ge2$. Donner le nombre de points du contour, puis pour $n=10$.
\label{fig-ex-8-11}\figureGeo{lit-11}
\end{exercice}\end{blocLivre}
\begin{blocLivre}{ex-8-12}\begin{exercice}[3]{Un tarif}
Une location coûte 8 euros fixes puis 3 euros par heure entière. Donner le prix pour $h$ heures, puis pour sept heures.
\end{exercice}\end{blocLivre}
\begin{blocLivre}{ex-8-13}\begin{exercice}[3]{Deux calculs à tester}
Pour $x=2$, comparer $x^2+4$ et $4x$. Peut-on conclure qu’ils sont toujours égaux ?
\end{exercice}\end{blocLivre}
\begin{blocLivre}{ex-8-14}\begin{exercice}[3]{Formule d’une suite}
Une suite de figures contient 5, 9, 13 puis 17 jetons. On ajoute toujours quatre jetons. Donner une formule pour l’étape $n$, la première étant $n=1$, puis le nombre à l’étape 25.
\end{exercice}\end{blocLivre}
\begin{blocLivre}{ex-8-15}\begin{exercice}[3]{Vérifier une vraisemblance}
Pour un rectangle de côtés $x+2$ et 3 cm, une formule proposée pour le périmètre est $P=3(x+2)$. Tester pour $x=4$ et corriger.
\label{fig-ex-8-15}\figureGeo{lit-15}
\end{exercice}\end{blocLivre}

\clearpage
